\documentclass[11pt]{article}
\usepackage{amsmath,hyperref}
\title{A resource handoff receipt in a warehouse and a railway model}
\author{}
\date{}
\begin{document}
\maketitle

A receiver that consumes a completed handoff can use the same rule for a warehouse delivery and a railway arrival. In the two source replays, however, the evidence comes from different events. A warehouse carrier can deliver as another carrier leaves its cell. A train can arrive and be removed in the same native tick, so its position afterward can say that it never occupied the resource. An accepted movement check can also leave a train where it was. Neither final occupancy nor an accepted check can replace an actual event.

\paragraph{What the receiver requires.}
Freeze one function on a single trusted transition. Its inputs state the resource's owner before application, the owner at application before optional terminal removal, the new task's service evidence at that resource, and the old task's remains-owed fact. It requires that the named old and new owners differ; that only the old owned this exact resource beforehand; that only the new owns it at application and produces the named arrival or delivery in the same tick; and that the named old task remains owed. A resource may have a more precise identity than a geometric cell. The receiver reads occurrence evidence after the transition; it neither authorizes the transition nor predicts service from promises. Its conclusion is only for the protected resource and the specified service event, not continuous occupancy between discrete cuts, indefinite availability, or settlement of the old task. The old task's type is a parameter and no equivalence of train trips and shelf returns is assumed. The fixture admits at most one scored service event per actor and resource at a cut, and one old task of the selected type per actor. Multiple concurrent jobs of the same type would require carrying a job identifier rather than just these keys.

\paragraph{Warehouse source.}
Use the \href{https://github.com/eternal-november/eos-research-programmes/pull/428}{selected native rotation harness} at RWARE 96fbc64. The adapter reads actual old/new cell ownership and native half reward plus the new delivery flag. The old carrier's still held shelf and prior delivery flag attest its separate return. Across the four replayed conditions, only the full admissible rotation yields the handoff. Suppressing a remote move or changing a carried shelf into a standing one prevents delivery, although an independent move still succeeds. For this discrete plant the application cut is also the post-step cell read. This adapter does not use any simulated grant to synthesize a movement or a delivery.

\paragraph{Railway source.}
Use \href{https://github.com/flatland-association/flatland-rl/tree/4856ebfe268d092d787d6b8536791a4fddb30772}{Flatland at 4856ebf}. Execute the selected native entry-candidate method, \texttt{MotionCheck}, result distribution, first-target bookkeeping, entry-point normalization, and grid-resource map. The adapter receives fixed on-map approach geometry, action validity and cell-exit status; native candidate and commit code decide which entry point is applied. It uses a stub speed-counter sink and a neutral stub for rewards. Thus the published motion and target effects are executed directly, but speed/rail generation, validity dynamics, a complete environment, and simulation performance are not evaluated.

An old train can leave the station block just as a following train claims it. The old train still has its own target ahead. If the follower is configured for removal, the executed native target handler first emits its target entry point and arrival tick, then removes it from the map. The native recorded target therefore reconstructs the application resource even when the map afterward is empty. Target identity by itself is not a fresh event: the replay also invokes the target handler at the following tick and verifies that the earlier recorded arrival is not refreshed. The same unchanged receiver works whether or not the follower is removed. Replacing application ownership with post-step map occupancy fails in the removal control.

Two rejected controls conserve the old right. The leading train may still be partway through the block, or an admitted malfunction can prevent the exit. Its positive \texttt{check\_resource} result means that keeping its current resource is allowed. The native conflict checker then denies the follower's attempt to acquire that same block, while an independent train still moves. Another control gives the map a grade-separated crossing. A horizontal arrival uses a different exclusive resource from the prior vertical resource even at the same coordinate. The horizontal arrival is real and does not prove a handoff of the vertical resource, including when the vertical train really leaves.

\paragraph{Checks and what transferred.}
The same receiver file, fixed before either adapter replay, is used for 12 source-level rail cases and four warehouse cases. The two positive rail cases exercise both arrival/removal outcomes; the warehouse positive exercises a simultaneous rotation. Controls block confirmation when old occupancy persists, the new arrival is absent, or the actual resource is different. The direct inspection of typed native movement and target outcomes gives the same answer, and no timing, operational advantage, or reduction in total proof effort is claimed. The portable part is the postcondition and its decision code. It costs separate source work to recover application ownership and live old debt. In particular this result does not transfer a movement policy, show that any source can produce the required evidence at an earlier decision cut, discharge R2 or I1, or qualify as an independently selected mission test.
\end{document}
